%! Potential Errors When Performing Tests — Unit 6, Lesson 6.7 — Warm-Up
\documentclass[10pt]{article}
\usepackage{apstats-article}
\usepackage{apstats-boxes}
\newcommand{\TallMath}[1]{$\displaystyle #1\rule[-1.4em]{0pt}{3.2em}$}

\begin{document}

\pageheader{Unit 6, Lesson 6.7}{Warm-Up}
\namedateperiod

\vspace{0.15in}

An environmental agency tests whether mercury contamination in a lake's fish is above the
safe limit of 5\% ($H_0: p = 0.05$ vs.\ $H_a: p > 0.05$, $\alpha = 0.05$).

\begin{enumerate}
  \item After testing 300 randomly selected fish, the agency gets $p\text{-value} = 0.038$ and
        \emph{rejects} $H_0$. However, the true contamination rate is actually 5\%.
        \begin{enumerate}[label=\alph*.]
          \item What kind of error did the agency make? \blank{3cm}
          \item Describe this error in context. \writelines{2}
          \item What is the probability of this type of error? \blank{2cm}
        \end{enumerate}

  \item Alternatively, suppose the true contamination rate is actually 9\%, but a different
        sample of 300 fish yields $p\text{-value} = 0.07$, so the agency \emph{fails to reject} $H_0$.
        \begin{enumerate}[label=\alph*.]
          \item What kind of error occurred? \blank{3cm}
          \item Describe this error in context. \writelines{2}
          \item In this scenario, which error seems more serious and why? \writelines{2}
        \end{enumerate}
\end{enumerate}

\end{document}
